\section{Fájl jogosultságok, ACL, setuid/setgid és sticky bitek}

Ez a fejezet részletesen tárgyalja a UNIX jogrendszer elemeit és finomhangolási lehetőségeit.

\subsection{chmod numerikus és szimbolikus mód}
\begin{lstlisting}[language=bash]
# Numeric example: owner rwx, group r-x, other ---
chmod 750 /opt/app
# Symbolic mode: add read permission to the group
chmod g+r file.txt
\end{lstlisting}

\subsection{setuid, setgid, sticky}
\begin{itemize}
  \item setuid (4000): futtatáskor a folyamat átveszi a fájl tulajdonosának jogosultságait (pl. passwd)
  \item setgid (2000): futtatáskor a csoport jogosultságot örökli; könyvtár esetén új fájlok csoportját beállítja
  \item sticky (1000): könyvtárban a fájlokat csak a tulajdonos vagy root törölheti (pl. /tmp)
\end{itemize}

\begin{lstlisting}[language=bash]
# setuid, setgid example
chmod 4755 /usr/bin/someprog
chmod 2755 /srv/shared
# set sticky bit on a directory
chmod +t /srv/shared
\end{lstlisting}

\subsection{ACL – finomhangolás}
\begin{lstlisting}[language=bash]
# Read ACL
getfacl /path/file
# Add ACL for user "alice" with rwx
setfacl -m u:alice:rwx /path/file
# Recursive ACL
setfacl -R -m u:alice:rwx /path/dir
\end{lstlisting}

Megjegyzés: ACL használatakor dokumentáljuk a szabályzatot, hogy elkerüljük a jogosultságok keveredését és auditálási problémákat.
